\uuid{b3Gt}
\exo7id{4468}
\titre{exo7 4468}
\auteur{quercia}
\organisation{exo7}
\datecreate{2010-03-14}
\isIndication{false}
\isCorrection{true}
\chapitre{Série numérique}
\sousChapitre{Autre}
\cpge{mpsi-mp2i.procedes-sommatoires-discrets.a, pcsi.series-numeriques.a, ptsi.series-numeriques.a, mpsi-mp2i.procedes-sommatoires-discrets.b, pcsi.series-numeriques.b, ptsi.series-numeriques.b, bcpst1.methodes-calcul, mp-mpi.series-entieres.b, pc.series-entieres.d, psi.series-entieres.d, pt.series-entieres.d, bcpst2.series-reelles}
\module{Analyse}
\niveau{L1}
\difficulte{}

\contenu{
\question{
Soient $a \in {[0,1[}$. \'Ecrire $\frac1{(1-a)^2}$ comme produit de deux
séries. En déduire la somme de la série $\sum_{k=0}^\infty ka^k$.
Calculer par la même méthode $\sum_{k=0}^\infty k^2a^k$.
}
\reponse{
$\frac{a}{(1-a)^2}$ et $\frac{a+a^2}{(1-a)^3}$.
}
}
